% D.86 candidate, unregistered. Root/independent full audits are pending.
\subsection{Infinite compact five-torsion in an actual cusp fibre}

\begin{statement}{D.86}{Open}{Candidate: an infinite torsion compact kernel for a cusp parameter}
Let
\[
 a=\sum_{q\ne5}a_q\kappa(q),\qquad a_q\longrightarrow0,
 \qquad \overline{\langle a_q:q\ne5\rangle}=\mathbb Z_5,
 \qquad
 \lambda=\sum_{q\ne5}a_q\frac{\log(q^4)}{4\log6},
 \qquad v_5(\lambda)=1.                                  \tag{D.86.1}
\]
Thus the five-valuation of \(a\) is zero, its global away-five
valuation image is primitive, and its local unit class is
\(\lambda\kappa(6)\). Arbitrary infinite \(c_0\) support is allowed.
For the actual Alexander module \(A_s\) of the pair \(s=(a,0)\),
the compact groups of D.63 satisfy
\[
 \begin{split}
 \mathcal S_s&\simeq\prod_{n\ge0}\mathbb F_5,\\
 \mathcal T_s&\simeq\prod_{n\ge0}\mathbb Z_5,\qquad
 L_5\simeq\mathcal T_s\times\mathbb Z_5^2,\\
 \mathcal K_s&\simeq
       \left(\prod_{n\ge0}\mathbb F_5\right)\times\mathbb Z_5^2.
 \end{split}                                             \tag{D.86.2}
\]
All isomorphisms are of compact topological groups. The map
\(L_5\to\mathcal T_s\) is the identity on the zero-parameter
boundary factor and kills the other two factors. The last
splitting is continuous and noncanonical.

Every pair \((a,0)\) here has an actual cusp section. In particular
\(a=\kappa(26)\) satisfies (D.86.1); its nonempty metabelian
section fibre is a torsor under
\(\prod_{\mathbb N}\mathbb Z_5\), by D.74. The infinite torsion
in (D.86.2) concerns the compact second-cohomology kernel, not
torsion in that first-cohomology translation group. Each finite
collection of the constructed nonzero dual classes occurs at
one finite coefficient level and is locally trivial there at
\emph{every} place. No single finite level is asserted to contain
all infinitely many independent classes.

This settles the stated globally dependent family only. It
makes no classification of all dependent pairs, no arbitrary
infinitely ramified first-cohomology rank computation, and no
first higher-derived obstruction computation.
\end{statement}

\begin{proof}[Candidate proof, pending full independent audits]
Write
\[
 \Lambda=\mathbb Z_5[[u,v]],\quad t_1=1+u,\quad t_2=1+v,
 \quad w=t_1t_2-1,\quad P=uvw,\quad D=A_s^D.
\]
A superscript \(D\) means Cartier dual; a star means ordinary
Pontryagin dual after the Cartier twist has been accounted for.
Let \(I_0\) be the actual compact untwisted operator image of
D.78, including its cyclotomic multiplier. Put
\(N_0=\ker(\chi:I_0\to\Gamma)\), with
\(\Gamma=\mathbb Z_5^\times\). The actual local surjectivity
\(G_{\mathbb Q_5}\to I_0\), the saturation of \(\kappa(6)\)
on its kernel, and the primitive cubic element of D.77--D.78
are used with their already proved integral meanings.

\paragraph{Actual global and local images.}
Let \(H=\mathbb Z_5\) act by multiplication by \(t_1^h\).
The untwisted action is unramified outside five. Actual inertia
at the primes in (D.86.1) therefore generates the whole pure
translation subgroup \(H\). The uniformly convergent cocycle
and local gauge construction of D.80/D.83 gives actual equalities
\[
              I_s=H\rtimes I_0,
       \qquad I_5=5H\rtimes I_0.                       \tag{D.86.3}
\]
The second equality uses \(\lambda\mathbb Z_5=5\mathbb Z_5\),
and both use the same zero-translation complement. The local
gauge lies in \(H\), so preserves the first equality. These
are faithful compact operator images, not formal semidirect
products assigned independently of Galois.

Put
\[
 q(u)=(1+u)^5-1,
 \qquad S(u)=q(u)/u=5+10u+10u^2+5u^3+u^4,
 \qquad B=\Lambda/(u),\quad R=\Lambda/(q),\quad J=uR.
                                                               \tag{D.86.4}
\]
The continuous translation coinduction of D.81/D.84 gives
natural identifications
\[
 H^1(I_s,D)=H^1(I_0,B^D),\qquad
 H^1(I_5,D)=H^1(I_0,R^D).                               \tag{D.86.5}
\]
Their restriction map is induced by the actual inclusion
\(B^D\subset R^D\). In particular its \emph{whole} kernel is
\[
 \mathcal R=(J^D)^{I_0}/\operatorname{im}((R^D)^{I_0}). \tag{D.86.6}
\]
The finite affine Chebotarev descent of D.78/D.83, at every
finite coefficient level before taking a union, gives
\[
                         \mathcal S_s^*\lhook\joinrel\longrightarrow
                         \mathcal R.                  \tag{D.86.7}
\]
This gives an upper bound; we will separately construct
infinitely many classes satisfying the common finite-level
local-triviality condition defining \(\mathcal S_s^*\).

\paragraph{The exact compact connecting module.}
Every \(g\in N_0\) has \(\mathcal B_g-1\in P\Lambda\).
There is an actual \(\delta\in N_0\) with
\(\mathcal B_\delta-1=P E\), \(E\in\Lambda^\times\),
by the local raw cubic primitivity of D.77. Consequently
\[
 (J^D)^{I_0}=(T(-1)_\Gamma)^*,\qquad T=J/PJ.            \tag{D.86.8}
\]
Here \((-1)\) means that the compact operator is
\(\chi^{-1}\rho_0\), exactly one Cartier untwist. The
\(\Gamma\)-action is induced by actual lifts in \(I_0\);
differences between lifts have already been killed by \(PJ\).
Compact coinvariants mean the quotient by the closed span of
operator differences.

Multiplication by the actual generator \(u\) identifies the
scalar module underlying \(J\) with
\[
             A=\Lambda/(S(u))=\mathcal O[[v]],
             \qquad \mathcal O=\mathbb Z_5[\zeta_5],
             \quad u=\zeta_5-1.                       \tag{D.86.9}
\]
The last equality follows from the Eisenstein polynomial in
D.77. Thus \(A\) is a domain, \(u\) is a uniformizer in
\(\mathcal O\), and \(A/(u)=\mathbb F_5[[v]]\).
In scalar coordinates \(T=A/(uvw)\), but its elements still
have the extra actual factor \(u\). We retain that factor in
all action calculations.

There is an exact sequence of actual compact modules
\[
 0\longrightarrow E\longrightarrow T\longrightarrow Q
       \longrightarrow0,
 \quad E=vwA/(uvwA)\simeq\mathbb F_5[[v]],
 \quad Q=A/(vw).                                      \tag{D.86.10}
\]
The generator of \(E\) in the \emph{actual} ideal is \(P=uvw\),
not the scalar \(vw\). The isomorphism for \(E\) uses
cancellation of the nonzero divisor \(vw\) in \(A\).
The ideals \((v)\) and \((w)\) have intersection \((vw)\):
modulo \(v\), \(w=u\) is a nonzero divisor in \(\mathcal O\),
and the converse divisibility follows by cancellation.
Therefore
\[
 0\longrightarrow Q\longrightarrow A/(v)\oplus A/(w)
       \longrightarrow A/(v,w)=\mathbb F_5\longrightarrow0.
                                                               \tag{D.86.11}
\]
Each term here carries the actual factor \(u\). In particular
\(Q\) is finite free of rank eight over \(\mathbb Z_5\).

\paragraph{The infinite torsion layer and the finite boundary.}
Let \(\gamma\in\Gamma\) and write
\(S_\gamma(t)=(t^\gamma-1)/(t-1)\). The actual multiplier
of \(\rho_0\) is
\(S_\gamma(t_1)S_\gamma(t_2)\mathcal B_g\).
On the generator \(P\), substitution contributes the additional
factor
\(S_\gamma(t_1)S_\gamma(t_2)S_\gamma(t_1t_2)\).
On \(E\), set \(u=0\) and reduce modulo five. Its compact
Cartier-untwisted action is consequently
\[
        f(v)\longmapsto
        \overline\gamma S_\gamma(1+v)^3
                   f((1+v)^\gamma-1).                \tag{D.86.12}
\]
There is no remaining beta term on this axis. Multiplication
by \(v^3\) identifies (D.86.12) with the ideal
\(v^3\mathbb F_5[[v]]\) inside the completed group ring
\(C=\mathbb F_5[[\mathbb Z_5]]\), with action
\[
                    t^x\longmapsto
                       \overline\gamma t^{\gamma x}.  \tag{D.86.13}
\]
The quotient \(C/v^3C\) has successive tame characters
\(\overline\chi,\overline\chi^2,\overline\chi^3\).
All are nontrivial on the order-four subgroup. Exact averaging
on that subgroup therefore identifies the coinvariants of
\(v^3C\) and \(C\).

The discrete ordinary dual of \(C\) is the continuous
\(\mathbb F_5\)-valued function space on \(\mathbb Z_5\).
Invariant functionals for (D.86.13) are functions satisfying
\[
 f(\gamma x)=\overline\gamma^{-1}f(x).
                                                               \tag{D.86.14}
\]
They vanish at zero. On each shell \(5^n\mathbb Z_5^\times\)
they are determined freely by their value at \(5^n\).
Continuity at zero forces those values to vanish for all
sufficiently large \(n\); conversely every finitely supported
sequence gives a continuous function satisfying (D.86.14).
It follows, by duality, that
\[
                         E(-1)_\Gamma
                               \simeq\prod_{n\ge0}\mathbb F_5.
                                                               \tag{D.86.15}
\]
This is the product compact topology, not a direct-sum topology.

For the two boundary modules in (D.86.11), after the Cartier
untwist the exact multiplier is respectively
\(S_\gamma(t)^2\) and
\(S_\gamma(t)^2t^{(1-\gamma)/2}\), on a scalar element
with actual generator \(u\). The second formula is the exact
reflection formula of D.39. Multiplication by \(u^2\), and
on the second boundary also by \(t^{-1/2}\), identifies both
modules with
\[
 I_5^2\subset\mathbb Z_5[\mathbb Z/5],
 \qquad t\longmapsto t^\gamma.                         \tag{D.86.16}
\]
The integral injectivity and equivariance are the same
cancellation argument as D.84.15, with the stated half-integral
power legitimate because two is a unit in \(\mathbb Z_5\).
The cokernel in (D.86.11), retaining the actual generator \(u\),
has character \(\overline\chi^2\) after untwisting. Its
order-four invariants and coinvariants vanish. Thus
\[
                      Q(-1)_\Gamma\simeq\mathbb Z_5^2.
                                                               \tag{D.86.17}
\]
Moreover the action on \(Q(-1)\) factors through the finite
order-four quotient: this is true on both finite group-ring
modules in (D.86.16), and (D.86.11) is an actual equivariant
injection.

Choose an actual lift \(k\) of \(2\in\Gamma\). Two
topologically generates \(\Gamma\), so all the compact
coinvariants above are cokernels of \(k-1\). The snake lemma
for (D.86.10) shows that the kernel of
\(E(-1)_\Gamma\to T(-1)_\Gamma\) is an image of
\(Q(-1)^{k=1}\). This last module is finite free over
\(\mathbb Z_5\), so its image in the exponent-five module
(D.86.15) is finite. A quotient of
\(\prod_{\mathbb N}\mathbb F_5\) by a finite subgroup is
again such a product: dualizing removes only a finite
codimension from a countably infinite-dimensional vector space.
Consequently (D.86.10), (D.86.15), and (D.86.17) give
\[
 T(-1)_\Gamma\simeq
             \left(\prod_{\mathbb N}\mathbb F_5\right)
                           \times\mathbb Z_5^2.        \tag{D.86.18}
\]
The split is continuous: choose lifts of the two basis vectors
of the free compact quotient, then use continuous
\(\mathbb Z_5\)-multiplication. This argument keeps the
possible finite snake-map image; it does not assume that
\(E\)-coinvariants inject.

\paragraph{The entire restriction kernel has exponent five.}
D.77 computes the actual invariant groups and their natural
inclusions. Globally only the \(b=0\) axis contributes; locally
the \(a\) and \(a-b\) axes each contribute one additional
copy of \(\mathbb Q_5/\mathbb Z_5\). Thus
\[
 (R^D)^{I_0}/(B^D)^{I_0}
                  \simeq(\mathbb Q_5/\mathbb Z_5)^2.   \tag{D.86.19}
\]
Its map into \((J^D)^{I_0}\) is injective by exactness.
By (D.86.8) and (D.86.18), the latter group is
\[
             \bigoplus_{\mathbb N}\mathbb F_5
                         \ \oplus\
                     (\mathbb Q_5/\mathbb Z_5)^2.
\]
Every divisible subgroup lies in its maximal divisible
subgroup. The image in (D.86.19) fills that subgroup:
dualizing its injection gives a surjection between free
\(\mathbb Z_5\)-modules of rank two, necessarily an
isomorphism. Therefore the actual quotient (D.86.6) satisfies
\[
                    \mathcal R\simeq
                         \bigoplus_{\mathbb N}\mathbb F_5.
                                                               \tag{D.86.20}
\]
In particular (D.86.7) implies that \(\mathcal S_s^*\)
has exponent five and at most countable dimension. This
upper calculation alone does not prove that its dimension
is infinite; the next steps address that separately.

\paragraph{Infinitely many actual global cocycles.}
The quotient \(J/uJ\), with its actual generator \(u\), is
\(\mathbb F_5[[v]]\). Its untwisted compact action is
\[
 f(v)\longmapsto
       \overline\gamma S_\gamma(1+v)
                      f((1+v)^\gamma-1).              \tag{D.86.21}
\]
All elements of \(N_0\) act trivially on this quotient.
Multiplication by \(v\) puts (D.86.21) inside the same
module (D.86.13); its missing constant quotient has
nontrivial character \(\overline\chi\). The shell argument
therefore gives an infinite-dimensional countable vector space
\[
 W:=((J/uJ)^D)^{I_0}\simeq\bigoplus_{\mathbb N}\mathbb F_5.
                                                               \tag{D.86.22}
\]
View \(W\) as a subspace of \((J^D)^{I_0}\).
Its intersection with the image of (D.86.19) has dimension
at most two, because it consists of elements of order five
in a divisible group of corank two. Hence its image in
\(\mathcal R\) is infinite-dimensional.

For \(\varphi\in W\), choose a concrete order-five extension
\(F\in R^D\) as follows. In
\[
                  R/5R=\mathbb F_5[[u,v]]/(u^5),
\]
let \(F\) apply the functional defining \(\varphi\) to the
coefficient of \(u^1\). Thus \(F|_J=\varphi\); it depends
on finitely many coefficients in \(v\), and on no powers
of \(u\) other than the first. Define
\[
 z_\varphi(g_0)=g_0F-F\in B^D\quad(g_0\in I_0),
 \qquad z_\varphi(h,g_0)=z_\varphi(g_0)\quad((h,g_0)\in I_s).
                                                               \tag{D.86.23}
\]
The first inclusion holds by invariance of \(\varphi\).
The second formula is an actual cocycle because \(H\)
acts trivially on \(B^D\). Its class is the connecting
image of \(\varphi\), under (D.86.5). Inflation from
\(I_s\) to \(G_{\mathbb Q}\) is injective in degree one
for these coefficients. Thus the infinite-dimensional image
just found consists of genuine nonzero global cohomology
classes. All cochains in (D.86.23) are continuous and have
values in a finite coefficient module, because \(F\) has
finite coefficient support and the actual action preserves
the augmentation filtration.

\paragraph{A bound for every unramified Frobenius primitive.}
Put \(\mathfrak a=(u,v)\), and let
\[
              M_n=(A_s/(5,\mathfrak a^n)A_s)^D.
\]
Choose \(n\ge2\) such that \(F\) and all values of
\(z_\varphi\) lie in \(M_n\). For any actual element
\(g=(h,g_0)\in I_s\), a primitive of
\(z_\varphi(g)\) for the cyclic group generated by \(g\)
exists in the single module \(M_{n+1}\). Here is a uniform
proof with all three cases distinguished.

If \(h\equiv0\pmod5\), the first-\(u\) functional \(F\)
is fixed by multiplication by \(t_1^h\): modulo five this
power has no term of degree one through four in \(u\).
Since \(H\) also fixes \(B^D\), (D.86.23) gives
\(gF-F=z_\varphi(g)\).

If \(\chi(g)\not\equiv1\pmod5\), integral conjugation
by an element of \(H\) removes the coordinate \(h\),
since \(\chi(g)-1\) is a unit. The cocycle is zero on
\(H\), and \(H\) fixes its values. Transporting the
primitive \(F\) from \((0,g_0)\) back by this conjugation
gives a primitive for \(g\) in \(M_n\). The conjugating
translation preserves this finite module. This is a
cocycle conjugation calculation, not an assumption that
the raw Frobenius translation was zero.

It remains that \(h\not\equiv0\pmod5\) and
\(\chi(g)\equiv1\pmod5\). On the compact coefficient
space modulo five set
\(U_g=\chi(g)^{-1}\rho_s(g)\). The exact low-degree
Fox/Ihara formulas of D.31/D.35 imply
\[
 U_g-1\text{ raises augmentation degree by at least one,}
 \qquad \operatorname{gr}(U_g-1)=\overline h\,u.
                                                               \tag{D.86.24}
\]
Indeed \(t_1^h=1+\overline h u+O(\mathfrak a^2)\);
\(S_\chi(t_i)\equiv1\pmod{\mathfrak a^2,5}\) when
\(\chi\equiv1\pmod5\); beta has no linear term;
and substitution \(t_i\mapsto t_i^\chi\) differs from
identity only in degree at least five on each variable.
The Cartier scalar is one modulo five. Multiplication by
\(\overline h u\) is injective on every homogeneous
polynomial piece of \(\mathbb F_5[u,v]\). Comparing the
least nonzero degree therefore proves that
\[
 U_g-1:\mathbb F_5[[u,v]]/\mathfrak a^n
          \longrightarrow
          \mathbb F_5[[u,v]]/\mathfrak a^{n+1}
                    \quad\hbox{is injective}.          \tag{D.86.25}
\]
Duality makes the actual \(g-1\) map
\(M_{n+1}\to M_n\) surjective. More explicitly, its
matrix differs from the transpose dual of \(U_g-1\)
by the invertible factor \(-U_g^{-T}\). Hence it has
a preimage of \(z_\varphi(g)\), with the same bound
\(n+1\) for every such \(g\). This proves the uniform
cyclic primitive assertion.

\paragraph{One finite global level for every place.}
Choose \(b\) with \(5^b\ge n+1\). By the \(c_0\)
condition, the set
\[
          E_n=\{q\ne5:a_q\notin5^b\mathbb Z_5\}
                                                               \tag{D.86.26}
\]
is finite. At every \(q\notin E_n\cup\{5\}\), the
module \(M_{n+1}\) is unramified: its inertia multiplier
is \(t_1^{a_q}\equiv1\pmod{5,\mathfrak a^{n+1}}\).
The cocycle (D.86.23) is also unramified there, because
it is inflated from the untwisted image \(I_0\).
Apply the preceding primitive construction to the actual
Frobenius element. It gives local triviality in the same
finite module \(M_{n+1}\) at all these infinitely many
primes.

At five, every element of the translation subgroup \(5H\)
fixes \(F\). Equation (D.86.23) is consequently the
coboundary of \(F\) on the actual whole group
\(I_5=5H\rtimes I_0\), already in \(M_n\).
At every \(q\in E_n\), the abelian inertia is nonzero.
D.43 and D.80's finite local Tate-duality/Mittag--Leffler
argument give
\(H^1(G_{\mathbb Q_q},D)=0\). Thus the local class
becomes a coboundary in some finite coefficient module
\((A_s/\mathfrak m^{r_q}A_s)^D\), where
\(\mathfrak m=(5,u,v)\). This is used for the finite
exceptional set only. Enlarge to one level containing
\(M_{n+1}\) and all those finitely many primitives.
The earlier primitives remain valid after inclusion, even
when this larger module has additional ramification.
Odd-primary cohomology at the real place vanishes.

We have proved that each class from (D.86.23) has a
representative in one finite global coefficient module
that is locally trivial at every place. Taking maxima
gives the same assertion for any finite collection.
No interchange of a filtered union with an infinite
product of local cohomology groups has occurred.
Their infinite-dimensional image therefore lies in
\(\mathcal S_s^*\). Combined with the exponent-five,
countable-dimensional upper bound (D.86.20), this gives
\[
                     \mathcal S_s^*\simeq
                            \bigoplus_{\mathbb N}\mathbb F_5.
\]
Pontryagin duality proves the first assertion of (D.86.2)
with its stated compact product topology.

\paragraph{The remaining factors and an actual cusp example.}
The nonzero classes \(a\) and \(a-b=a\) have global
depth zero and local depth one; the zero class \(b\)
has infinite depth both globally and locally. D.77's
natural boundary description therefore gives exactly
\(L_5\to\mathcal T_s\) and its kernel \(\mathbb Z_5^2\)
as stated. Its continuously split compact sequence gives
\(\mathcal K_s\simeq\mathcal S_s\times\mathbb Z_5^2\).
This concerns only D.63's compact separated kernel and
asserts no splitting of the unrestricted nonseparated
cochain quotient.

The cusp at zero has inertia parameter \(a\) and exact
abelian pair \((a,0)\); its inertia cocycle supplies an
actual geometrically pro-five section. One can also take
zero parameters at all other coefficient primes to get
an actual full profinite cusp section.
For \(a=\kappa(26)=\kappa(2)+\kappa(13)\), the away
coefficients are primitive and the five-valuation is zero.
Since \(26=1+25\),
\[
                 v_5(\log26)=2,\qquad v_5(\log6)=1,
                 \qquad v_5(\lambda)=1.
\]
Thus this explicit cusp section has the compact kernel
in (D.86.2). Its twist is finitely ramified, so the
already evaluated D.74 theorem computes its nonempty
metabelian fibre as stated, without identifying that
whole fibre with the cusp packet.
\end{proof}
